\documentclass[tikz,border=8pt]{standalone}
\usepackage{ctex}
\usetikzlibrary{arrows.meta,positioning,fit,calc,backgrounds}
% Shared visual language for LFS architecture diagrams.
\RequirePackage{../../mymath}

\definecolor{LFSBlue}{HTML}{2563EB}
\definecolor{LFSBlueSoft}{HTML}{DBEAFE}
\definecolor{LFSGreen}{HTML}{15803D}
\definecolor{LFSGreenSoft}{HTML}{DCFCE7}
\definecolor{LFSOrange}{HTML}{C2410C}
\definecolor{LFSOrangeSoft}{HTML}{FFEDD5}
\definecolor{LFSRed}{HTML}{B91C1C}
\definecolor{LFSRedSoft}{HTML}{FEE2E2}
\definecolor{LFSPurple}{HTML}{7E22CE}
\definecolor{LFSPurpleSoft}{HTML}{F3E8FF}
\definecolor{LFSGray}{HTML}{475569}
\definecolor{LFSGraySoft}{HTML}{F1F5F9}

\tikzset{
  lfs picture/.style={
    font=\small,
    node distance=8mm and 12mm,
    >=Latex,
    every path/.style={line cap=round, line join=round}
  },
  block/.style={
    draw=LFSBlue,
    fill=LFSBlueSoft,
    rounded corners=2pt,
    line width=0.8pt,
    align=center,
    inner xsep=7pt,
    inner ysep=5pt,
    minimum height=10mm
  },
  compute/.style={
    block,
    draw=LFSPurple,
    fill=LFSPurpleSoft
  },
  storage/.style={
    block,
    draw=LFSGreen,
    fill=LFSGreenSoft
  },
  interface/.style={
    block,
    draw=LFSOrange,
    fill=LFSOrangeSoft
  },
  risk/.style={
    block,
    draw=LFSRed,
    fill=LFSRedSoft
  },
  neutral/.style={
    block,
    draw=LFSGray,
    fill=LFSGraySoft
  },
  decision/.style={
    diamond,
    draw=LFSOrange,
    fill=LFSOrangeSoft,
    line width=0.8pt,
    align=center,
    inner sep=2pt,
    aspect=2
  },
  group/.style={
    draw=LFSGray,
    rounded corners=3pt,
    dashed,
    line width=0.7pt,
    inner sep=6pt
  },
  arrow/.style={
    -{Latex[length=2.4mm,width=1.6mm]},
    draw=LFSGray,
    line width=0.9pt
  },
  data arrow/.style={
    arrow,
    draw=LFSBlue
  },
  control arrow/.style={
    arrow,
    draw=LFSOrange,
    dashed
  },
  residual/.style={
    arrow,
    draw=LFSGreen,
    rounded corners=4pt
  },
  label text/.style={
    font=\scriptsize,
    text=LFSGray,
    fill=white,
    inner sep=1.5pt
  },
  title/.style={
    font=\bfseries,
    text=LFSGray,
    align=center
  }
}

\begin{document}
\begin{tikzpicture}[my picture,node distance=8mm and 10mm]

\node[neutral,text width=25mm] (guidance) {项目与全局指令\\\scriptsize \texttt{AGENTS.md}};
\node[compute,text width=23mm,right=of guidance] (merge) {作用域合并\\\scriptsize 形成适用指令 $I_t$};

\node[interface,text width=25mm,below=13mm of guidance] (inputs) {记忆候选来源\\\scriptsize 合格历史任务、人工更正};
\node[compute,text width=23mm,right=of inputs] (consolidate) {后台抽取与归并\\\scriptsize 脱敏、去重、冲突处理};
\node[storage,text width=23mm,right=of consolidate] (store) {本地记忆状态\\\scriptsize 摘要、条目、来源证据};
\node[compute,text width=23mm,right=of store] (retrieve) {按任务检索\\\scriptsize 形成召回记忆 $M_t$};
\node[block,text width=24mm,right=of retrieve] (context) {当轮任务上下文\\\scriptsize $q_t,I_t,M_t,O_t$};

\draw[data arrow] (guidance) -- (merge);
\draw[data arrow] (merge.east) -| (context.north);
\draw[data arrow] (inputs) -- (consolidate) -- (store) -- (retrieve) -- (context);

\node[neutral,text width=31mm,below=10mm of store] (controls) {读取与贡献控制\\\scriptsize 会话开关、全局配置};
\draw[control arrow] (controls.west) -| (consolidate.south);
\draw[control arrow] (controls.east) -| (retrieve.south);

\node[group,fit=(guidance)(merge),label={[title]above:稳定指令路径}] {};
\node[group,fit=(inputs)(consolidate)(store)(retrieve)(controls),label={[title]above:可召回记忆路径}] {};

\end{tikzpicture}
\end{document}
